\documentclass[a4paper,10pt,twoside]{article}
\usepackage{conf_template}
\usepackage[utf8]{inputenc}
\setcounter{section}{0}
\setcounter{figure}{0}
\setcounter{table}{0}
\setcounter{equation}{0}
\setcounter{secnumdepth}{1}
\usepackage{multicol}
\usepackage{mathtools}
\usepackage[utf8]{inputenc}
\setcounter{secnumdepth}{1}
\usepackage{enumitem}


\begin{document}
% ================================= Начало Вашей статьи. Всё, что выше - не изменять ! ================================
%

\authors {% авторы статьи, между инициалами - пробел
Т.~В.~Тиханович
К.~А.~Хаджинова
}%конец списка авторов
%
\topic{Использование Bash скриптов для работы с операционной системой Windows}
%
% ------------- Аннотация -----------------------
\annotation{%% начало текста аннотации
В статье рассматривается использование Bash-скриптов для работы с операционной системой Windows, что открывает новые возможности для автоматизации задач и управления системой.
}% конец текста аннотации
\begin{multicols}{2}
% =================================== Начало основного текста ===========================================================
% ===== Первый раздел статьи =================================================
%
\section*{Введение}

Bash -скрипты, чаще ассоциируются с операционной системой (ОС) Linux, однако они также могут быть использованы для работы с ОС Windows. В настоящее время ОС Windows является самой популярной ОС для персональных компьютеров и используется в мире бизнеса и домашних компьютеров благодаря своей удобной интерфейсной среде. Она обладает широким набором функций и возможностей, но для опытных пользователей, привыкших к управлению системой через командную строку, интуитивные интерфейсы могут казаться несколько ограничивающими. В таких случаях можно использовать Bash-скрипты, которые позволяют автоматизировать выполнение задач и выполнять сложные действия, которые не предусмотрены стандартными средствами Windows.
% конец текста первого раздела статьи
% ===== Второй раздел статьи ===================================================
%
\section{Что такое Bash скрипты?}
\content{% начало текста 2 раздела статьи
Bash скрипты – это текстовые файлы, содержащие команды Bash, которые выполняются последовательно. 
Bash или Bourne Again SHell – это командный интерпретатор командной строки, используемый в ОС семейства Linux и Unix, и позволяет автоматизировать задачи в ОС. Он также доступен в Windows 10 Anniversary Update и более поздних версиях. Хотя Windows использует Windows PowerShell в качестве своего основного командного интерпретатора, Bash-скрипты все равно могут быть полезными для работы с Windows.
}% конец текста второго раздела статьи
%
% ======-Третий раздел статьи ===================================================
%
\section{Установка Bash на Windows}
\content{% начало текста 3 раздела статьи
Использование Bash-скриптов для работы с ОС Windows возможно благодаря Windows Subsystem for Linux (WSL). WSL позволяет запускать полноценную систему Linux внутри Windows 10/11, включая командную оболочку Bash. Установка WSL и Bash на Windows 10/11 относительно проста и позволяет использовать команды Linux напрямую в Windows, без необходимости установки виртуальной машины или двойной загрузки. Для установки Bash на Windows 10/11 необходимо включить подсистему Windows для Linux и загрузить необходимый дистрибутив Linux, такой как Ubuntu, из Магазина Windows. После установки можно использовать Bash для запуска скриптов и выполнения команд Linux [1]. WSL предназначен в первую очередь для разработчиков, чтобы они могли использовать инструменты Linux напрямую на Windows, и является мощным инструментом для управления Unix и Linux серверами [2].
Современные версии Windows поддерживают Bash, что позволяет пользователям запускать Bash-скрипты напрямую на своих машинах. Для этого нужно установить подсистему Windows для Linux (WSL): установить Bash, активировав функцию "Windows Subsystem for Linux" в настройках Windows. Это можно сделать через "Параметры" -> "Обновление и безопасность" -> "Для разработчиков". Затем следует включить WSL и установить нужный дистрибутив, например, Ubuntu [3].
После установки Bash можно создавать и запускать Bash-скрипты на Windows. Это позволяет выполнять различные задачи, такие как создание, копирование, перемещение файлов и папок, управление службами Windows, выполнение команд PowerShell и многое другое.
Как написать Bash-скрипт?
После установки WSL пользователи получают доступ к командной строке Linux, где они могут создавать и выполнять Bash-скрипты. Для создания скрипта можно воспользоваться любым текстовым редактором, таким как Nano, Vim или даже Notepad++. В начале скрипта необходимо указать интерпретатор, который будет его выполнять. В случае Bash-скриптов – это будет строка вида:
 \image{[width=0.3\columnwidth]{pic.1.png}}
Далее можно вписывать команды Bash, которые будут выполняться последовательно. Каждая команда должна заканчиваться символом новой строки.
Пример простого Bash-скрипта (например, script.sh):
 \image{[width=0.8\columnwidth]{pic.2.png}}
Как запустить Bash скрипт?
Для запуска Bash-скриптов, сохраненных на Windows, достаточно выполнить команду bash script.sh в командной строке WSL :
 \image{[width=0.3\columnwidth]{pic.3.png}}
Где script.sh – имя скрипта.
Это позволяет использовать Bash-скрипты для автоматизации повседневных задач, таких как управление файлами и папками, обновление и установка программ [4].
}% конец текста третьего раздела статьи
% ===== Четвертый раздел статьи ===================================================
%
\section{Интеграция с системой Windows}
\content{ % начало текста 4 раздела статьи
Bash-скрипты могут взаимодействовать с системой Windows, вызывая исполняемые файлы, работая с реестром или даже управляя службами. Это открывает дополнительные возможности для создания мощных инструментов автоматизации.
Пример вызова исполняемого файла из Bash-скрипта:
 \image{[width=0.8\columnwidth]{pic.4.png}}
Например, можно написать Bash-скрипт, который копирует все файлы из одной папки в другую:
  \image{[width=0.8\columnwidth]{pic.5.png}}
   Кроме того, можно использовать Bash-скрипты для выполнения команд PowerShell на Windows. Например, можно написать Bash-скрипт, который запускает PowerShell скрипт для создания нового пользователя:
  \image{[width=8cm,height=2.5cm]{pic.6.png}}
Этот скрипт использует команду powershell для запуска PowerShell скрипта, который создает нового локального пользователя с указанным именем и паролем.
} 

% конец текста четвертого раздела статьи
%
% ===== Пятый раздел статьи ===================================================
%
\section{Примеры использования Bash-скриптов для работы с Windows}
\content{% начало текста 5 раздела статьи
Несколько примеров того, как можно использовать Bash-скрипты для работы с Windows:
\item1.	Автоматизация выполнения рутинных задач. Bash предоставляет широкий спектр команд и функций для работы с файлами, папками, текстовыми данными и другими ресурсами операционной системы. Можно создавать скрипты, которые выполняют серию операций, таких как запуск приложений, копирование, перемещение, удаление файлов, а также обработку текстовых файлов и выполнение других системных команд.
\item2.	Выполнение сложных действий, которые не предусмотрены стандартными средствами Windows. Например, можно написать скрипт, который будет устанавливать или удалять программное обеспечение, создавать или удалять пользователей или изменять системные настройки.
\item3.	Интеграция с другими инструментами и сценариями. Bash поддерживает вызов внешних программ и команд, что позволяет использовать его в сочетании с другими инструментами разработки и автоматизации. Например, можно написать Bash-скрипт, который запускает сборку вашего проекта, после чего вызывает другие инструменты для тестирования и развертывания.
\item4.	Использование Bash-скриптов в качестве утилит. Bash-скрипты можно использовать в качестве утилит, которые можно вызывать из командной строки. Например, можно написать скрипт, который будет выводить информацию о системе или выполнять диагностику.
\item5.	Использование при разработке и тестировании программного обеспечения под Windows. Можно написать скрипт, который автоматически компилирует и собирает ваше приложение, запускает набор тестов и генерирует отчеты о результатах. Это может значительно упростить процесс разработки и улучшить эффективность вашей работы.
Несмотря на все преимущества, стоит отметить, что использование Bash-скриптов на Windows требует некоторых знаний и понимания командной строки и основных концепций Linux. Также следует помнить о некоторых различиях в командах и синтаксисе между Windows и Linux, которые могут потребовать небольших изменений в скриптах.
}% конец текста пятого раздела статьи
% ===== Шестой раздел статьи ===================================================
%
\section{Преимущества Bash-скриптов перед «пакетными» скриптами (.bat/.cmd) или скриптами PowerShell (.ps1)}
\content{% начало текста 6 раздела статьи
\item1.	Портативность: Bash-скрипты являются стандартом в Unix-подобных системах, включая Linux. Это означает, что скрипты, написанные на Bash, могут быть запущены на различных платформах, включая Windows с помощью Windows Subsystem for Linux (WSL) [5].
\item2.	Простота и удобство: Bash-скрипты обычно имеют более простый и понятный синтаксис по сравнению с cmd или PowerShell. Они используют привычные команды Unix, такие как cd, ls, mkdir, что делает их более доступными для новичков.
\item3.	Мощные возможности: Bash-скрипты предоставляют широкий набор возможностей для автоматизации задач. Они поддерживают циклы, условные операторы, функции и многое другое. Вы можете использовать их для обработки файлов, управления процессами, работе с сетью и многого другого.
\item4.	Интеграция с другими инструментами: Bash-скрипты могут легко взаимодействовать с другими утилитами и программами в Unix-подобных системах. Вы можете использовать их в сочетании с командами grep, sed, awk и другими инструментами для более сложных операций обработки данных.
\item5.	Большое сообщество и документация: Bash имеет большое сообщество пользователей и разработчиков, что означает, что вы можете найти множество ресурсов, документации и примеров кода для изучения и использования.
Однако, несмотря на эти преимущества, cmd.exe и PowerShell также имеют свои сильные стороны:
-	Cmd.exe – интерпретатор командной строки в Windows и может быть полезным для выполнения простых задач. 
-	PowerShell, с другой стороны, предоставляет более мощные возможности для автоматизации и управления Windows-системами.
В конечном счете, выбор между Bash, cmd.exe и PowerShell зависит от потребностей и предпочтений пользователя.
}% конец текста шестого раздела статьи
%
\section{Выводы}
Использование Bash-скриптов на операционной системе Windows через подсистему Windows для Linux предоставляет удобный способ автоматизации решения задач, управления системой, разработку программного обеспечения и выполнения сложных действий через командную строку. Это особенно полезно для тех, кто привык к командной строке Linux и хочет использовать ее преимущества в среде Windows.
%
\section*{Использованные источники}
%\references{}
\ListReferences{% начало списка литературы
\item  How to Install Linux Bash Shell on Windows [Электронный ресурс].–Режим доступа: https://itsfoss.com/install-bash-on-windows/.~–vДата доступа: 23.01.2024.
\item  Bash on Windows: What Does It Mean? [Электронный ресурс].–Режим доступа: https://www.linux.com/news/bash-windows-what-does-it-mean/.~–~Дата доступа: 23.01.2024.
\item  How To Install Bash On Windows 10 [Электронный ресурс].–Режим доступа: https://hackernoon.com/how-to-install-bash-on-windows-10-lqb73yj3.~–~Дата доступа: 23.01.2024.
\item  How To Enable Bash for Windows 10 & 11 [Электронный ресурс].–Режим доступа: https://www.onlogic.com/company/io-hub/how-to-enable-bash-for-windows-10-and-11/.~-~Дата доступа: 23.01.2024.
\item  How To Use Bash Shell Natively On Windows 10 [Электронный ресурс].–Режим доступа: https://www.geeksforgeeks.org/use-bash-shell-natively-windows-10/.~–~ Дата доступа: 23.01.2024.
}% ---- конец списка литературы
%
\end{multicols}

% ----------- конец данных об авторах
%
\authorFIO{Тиханович Татьяна Викторовна}
\authorAbout{
зам.декана ФИТиУ,
старший преподаватель кафедры ИТАС, БГУИР,
tihanovich@bsuir.by.
}
\authorFIO{Хаджинова Ксения Александровна}
\authorAbout{
студентка группы 328506,
кафедра ИТАС, БГУИР, 
xju2005@gmail.com 
}
% ----------- конец данных об авторах
\end{document}